\documentclass[10pt]{article}
\providecommand{\FIGDIR}{fig}
\newcommand{\DIGESTDATE}{6 October 2026}
\input{preamble}

\begin{document}

% ==================================================================== MASTHEAD
\thispagestyle{empty}
\begin{tikzpicture}[remember picture, overlay]
  \fill[Deep] (current page.north west) rectangle
      ($(current page.north east)+(0,-67mm)$);
  \fill[Amber] ($(current page.north west)+(0,-67mm)$) rectangle
      ($(current page.north east)+(0,-68.4mm)$);
  \foreach \x/\y/\r in {22/12/0.9, 41/26/1.5, 63/9/0.7, 88/31/1.1, 112/17/1.9,
                        138/29/0.8, 159/13/1.3, 178/54/1.0, 196/21/1.6,
                        30/44/0.6, 74/57/1.2, 125/48/0.9, 168/6/0.8, 191/59/0.7,
                        52/18/1.0, 100/60/0.6, 148/20/1.1, 66/49/1.4, 205/45/0.8}
    \fill[Sky, opacity=0.30] ($(current page.north west)+(\x mm,-\y mm)$) circle (\r mm);
\end{tikzpicture}

\vspace*{4mm}
{\sffamily\fontsize{7.6}{9}\selectfont\color{Sky}%
 AUTOMATED DAILY LITERATURE DIGEST\;\textbullet\;PhD RESEARCH WATCH\par}
\vspace{2.2mm}
{\sffamily\bfseries\fontsize{27}{29}\selectfont\color{white}%
 Particulate Matter\par}
{\sffamily\fontsize{27}{29}\selectfont\color{Sky}%
 Monitoring \& Health Effects\par}

\vspace{5.0mm}
{\sffamily\fontsize{8.4}{10}\selectfont\color{white}%
 Issue 2026-10-06\;\;\textcolor{Amber}{\textbar}\;\;Entry window 6 October 2026 (full day)\;\;%
 \textcolor{Amber}{\textbar}\;\;11 records screened in scope\par}
\vspace{18mm}

\noindent
\begin{minipage}[t]{0.190\linewidth}\metric{11}{RECORDS IN SCOPE\\(61 EXCLUDED)}{Deep}\end{minipage}\hfill
\begin{minipage}[t]{0.190\linewidth}\metric{6}{OF 10 SUBTOPIC\\BANDS POPULATED}{Teal}\end{minipage}\hfill
\begin{minipage}[t]{0.190\linewidth}\metric{0}{EFFECT ESTIMATES\\WITH A CI}{Sky}\end{minipage}\hfill
\begin{minipage}[t]{0.190\linewidth}\metric{2}{TIER-A\\RECORDS}{Amber}\end{minipage}\hfill
\begin{minipage}[t]{0.190\linewidth}\metric{4}{RECORDS CARRIED ON\\METADATA ALONE}{Clay}\end{minipage}

\vspace{6mm}

\section*{\sffamily\bfseries\color{Deep}Signal of the day}
\vspace{-1.5mm}{\color{Amber}\rule{\linewidth}{1.1pt}}\vspace{2.5mm}

\begin{keybox}
{\sffamily\bfseries\small\color{Deep}An emission-control policy moved mean particle diameter by
27\,\% --- which is a calibration drift no sensor firmware will ever report.}\\[1.4mm]
\footnotesize
Dai and colleagues (\textit{ACP}) ran single-particle aerosol mass spectrometry through winter and summer
emission-control periods in Yangzhou against matched normal periods, defining sulfur-to-nitrogen ratios in
the gas phase, the particle phase and by particle number. Cutting NO$_x$ while SO$_2$ stayed flat raised
the gas-phase ratio, lifted the particle-phase ratio by \textbf{$\sim$57\,\%} across particle classes, and
expanded mean vacuum aerodynamic diameter by \textbf{27\,\%}. Enrichment was largest in particles carrying
both black and organic carbon, least in BC-free inorganic particles.
\textbf{Caveat:} control periods are confounded with season and meteorology, and ``causal inference'' is a
strong label for an observational series with co-varying controls.
\textbf{Why it matters for sensing:} a nephelometric node's mass estimate is a fixed assumption about size
distribution and refractive index wrapped around a scattering measurement. A 27\,\% shift in mean diameter
changes the scattering cross-section per unit mass directly --- so a network calibrated before an emission
control reads differently after it, with no hardware change and nothing in the diagnostics to show for it.
Co-location drift studies attribute this to the sensor. It is not the sensor.\par
\end{keybox}

\vspace{3mm}

\section*{\sffamily\bfseries\color{Deep}What changed today}
\vspace{-1.5mm}{\color{Amber}\rule{\linewidth}{1.1pt}}\vspace{2.5mm}

\footnotesize
\begin{itemize}

\item \textbf{Infiltration above 0.6 in naturally ventilated classrooms.} Ten Klang Valley primary schools
across five land-use classes: indoor 24\,h PM$_{2.5}$ peaked at mixed industrial/port schools
(\textbf{35.33\,$\pm$\,9.25}\,$\mu$g\,m$^{-3}$), outdoor at high-traffic schools (34.50\,$\pm$\,20.64), with
mean \textbf{F$_{\textrm{inf}}$ $>$ 0.6}. Two buildings per class confounds land use with building, and
F$_{\textrm{inf}}$ from indoor--outdoor ratios absorbs any indoor source --- but if it holds, outdoor
network data is a usable proxy for children's exposure in this climate.

\item \textbf{Brain structure, without numbers.} 4,276 adults with repeated MRI: higher long-term PM$_{2.5}$
and PM$_{10}$ associated with accelerated grey-matter loss and greater white-matter-hyperintensity
progression. The abstract reports \textit{no} effect size, interval, exposure contrast, cohort name or
country, so the record is logged and cannot be entered as an estimate.

\item \textbf{Italy's national attributable-burden platform.} PM$_{2.5}$ and NO$_2$ modelled at 1\,km by
random forest for 2016--2024, population-weighted to municipality level against the 2021 census, and run
through HRAPIE-2 functions to the WHO 5\,$\mu$g\,m$^{-3}$ counterfactual. The surface is validated against
the monitors that train it, and municipality-scale population weighting erases the within-city contrast ---
which is precisely the gap a dense network fills.

\item \textbf{Port outflow at 4--10\,$\mu$g\,m$^{-3}$.} Four seasons of HR-ToF-AMS at a coastal receptor
downwind of Busan: five PMF organic factors, frequent ultrafine growth events, and oxidised primary organic
aerosol reaching \textbf{36\,\%} in autumn under north-easterly port transport. Growth events at that mass
loading are invisible to a mass-reporting network.

\end{itemize}

\vspace{2mm}

\begin{keybox}
{\sffamily\bfseries\footnotesize\color{Deep}Window continuity}\\[1.2mm]
{\fontsize{7.2}{8.8}\selectfont
The 5 October issue closed with \texttt{last\_entry\_date} at 5 October. This issue collects
\textbf{6 October in full}, continuous with it, and is the last day covered by the cached harvest. Harvest
and screening were completed on the 6 October run; this run authored from that cache without re-querying.
\textbf{Dailies for 7 and 8 October have not yet been harvested} and the W40 weekly (27 September--3
October) is still owed; both follow.\par}
\end{keybox}

\vspace{3mm}
\par\vskip 0pt plus 18\baselineskip\penalty-200\vskip 0pt plus -18\baselineskip
\section*{\sffamily\bfseries\color{Deep}Corpus analytics}
\vspace{-1.5mm}{\color{Amber}\rule{\linewidth}{1.1pt}}\vspace{2.5mm}

\noindent\includegraphics[width=\linewidth]{\FIGDIR/f1_subtopics.png}
\figcap{Subtopic assignment of the 11 in-scope records. \textbf{Six of ten bands populated};
\textit{Sensing} leads with three, \textit{Burden}, \textit{Exposure assessment} and \textit{Mechanistic
toxicology} two each, and one each for \textit{Neuro} and \textit{Occupational \& indoor}. The measurement
side holds 8 of 11 records --- the highest share of the backfill.}

\vspace{2mm}
\noindent\includegraphics[width=\linewidth]{\FIGDIR/f2_design_endpoint.png}
\figcap{Left --- architecture: \textbf{four} measurement campaigns (the AMS receptor
site, the single-particle series, the school indoor--outdoor survey and the Kazakh elemental sampling),
\textbf{four} metadata-only deposits, and one each for the imaging cohort, the national health-impact
model and the coal-fire laboratory rig. Right --- only two records report a health endpoint, and one of
those (the brain-imaging cohort) reports it without a number.}

\vspace{2mm}
\noindent\includegraphics[width=\linewidth]{\FIGDIR/f3_metric_geography.png}
\figcap{Left --- particle metric: PM$_{2.5}$ only \textbf{four}, unspeciated PM or
emissions three, PM$_{2.5}$ + PM$_{10}$ jointly two, and one each for size-resolved number distribution
(the submicron AMS record) and optical properties (the HULIS deposit). Right --- geography: China and
global / not stated \textbf{three} each, and one each for Europe, Southeast Asia, Central Asia, East Asia
excluding China, and South Asia --- Kazakhstan is the first Central Asian record since 29 September.}

\vspace{2mm}
\noindent\includegraphics[width=\linewidth]{\FIGDIR/f6_lifecourse.png}
\figcap{Life-course windows: childhood one (the Klang Valley school measurement) and older
adults one (the repeated-MRI cohort). In utero, adolescence and working age have no record --- a day
dominated by atmospheric measurement rather than by population studies.}

\vspace{2mm}
\noindent\includegraphics[width=\linewidth]{\FIGDIR/f4_heatmap.png}
\figcap{Subtopic $\times$ study architecture for the 11 records. Measurement campaigns occupy
sensing, exposure and occupational; the metadata-only block spans sensing, exposure and mechanistic
toxicology; burden holds the only two modelling-type records, one national and one at bench scale.}

\vspace{2mm}

{\fontsize{7.4}{9}\selectfont\color{Slate}\textbf{No forest plot this issue.} No admitted
abstract reports a PM effect estimate with a confidence interval. The brain-imaging cohort is the only
record with an epidemiological estimand, and its abstract reports direction without magnitude, interval or
exposure contrast; the Kazakh hazard quotients and the Italian attributable fractions are risk-assessment
outputs, not estimated associations, and are excluded from the panel by the same rule that has excluded
them all series. The f5 panel is absent by design, not by build failure.\par}

\vspace{4mm}

\par\vskip 0pt plus 24\baselineskip\penalty-200\vskip 0pt plus -24\baselineskip
\section*{\sffamily\bfseries\color{Deep}Paper-level digest}
\vspace{-1.5mm}{\color{Amber}\rule{\linewidth}{1.1pt}}\vspace{2.5mm}

{\fontsize{7.4}{9}\selectfont\color{Slate}Ordered by cluster. Each entry gives the
methodological core and the finding that matters, not an abstract paraphrase. Records
marked \textit{metadata only} carry no recoverable abstract and assert no finding.\par}

\band{Neuro / mental health\hfill 1 record}{Deep}

\paperentry{Wang et al. 2026c}
{Long-Term Air Pollution Exposure and Brain Structural Decline: A Longitudinal Analysis}
{Environ Pollut}
{4,276 middle-aged and older adults with repeated brain MRI; global, grey- and white-matter volumes and white-matter hyperintensity burden regressed on modelled residential PM$_{2.5}$ and PM$_{10}$ in fully adjusted multivariable models. Higher long-term particulate exposure was associated with accelerated global and regional grey-matter loss and greater WMH progression, which the authors read as both a neurodegenerative and a cerebrovascular pathway. Limitation: the abstract gives no effect sizes, no confidence intervals, no exposure contrast and does not name the cohort or country - unusual for a quantitative imaging analysis, and it makes the result unassessable as reported; repeat-MRI subsamples are also selected for survival and compliance. Relevance: moderate - brain structure is the mediator between the UFP dementia-mortality signal and clinical disease, but this record cannot be entered as an effect estimate.}
{10.1016/j.envpol.2026.129269}

\par\vskip 0pt plus 10\baselineskip\penalty-200\vskip 0pt plus -10\baselineskip
\band{Mechanistic toxicology\hfill 2 records}{Violet}

\paperentry{Sharma et al. 2026}
{Characterization of carbonaceous particles generated from three flame modes using a single burner to enable properties-based exposure studies}
{Aerosol Sci Technol}
{\textsc{metadata only.} Metadata-only record: Taylor \& Francis deposited title, authors and DOI on 6 Oct with no abstract. By title, a single burner operated in three flame modes to generate carbonaceous particles of controlled and differing properties, explicitly to support properties-based rather than mass-based exposure studies. That framing is the open problem behind every toxicology record in this backfill: dose expressed as mass hides composition, and a generator that varies one property at a time is how it gets separated. Relevance: high; on the retry list.}
{10.1080/02786826.2026.2737177}

\paperentry{Yadav et al. 2026}
{The dual role of atmospheric HULIS in inducing climate forcing and oxidative stress in the northwestern Himalayas}
{Environ Pollut}
{\textsc{metadata only.} Metadata-only record: Elsevier deposited title, authors and DOI on 6 Oct with no abstract. By title, HULIS measured for both light absorption (brown-carbon forcing) and oxidative stress, in a high-altitude Himalayan setting. HULIS is the fraction that links the optical property a low-cost sensor responds to with the oxidative potential that toxicology implicates, so a single study reporting both is more useful than two reporting one each. Relevance: high; on the retry list.}
{10.1016/j.envpol.2026.129300}

\par\vskip 0pt plus 10\baselineskip\penalty-200\vskip 0pt plus -10\baselineskip
\band{Exposure assessment \& modelling\hfill 2 records}{Clay}

\paperentry{Mukhamediyarov et al. 2026}
{Elemental Composition, Sources, and Health Risks of Ambient PM$_{2.5}$ in Ust-Kamenogorsk (Oskemen) and Pavlodar, Kazakhstan}
{Atmosphere}
{PCA separates a non-ferrous metallurgical profile in Ust-Kamenogorsk from coal fly ash, ferrous metallurgy and petrochemicals in Pavlodar. The risk arithmetic then reports cumulative hazard indices exceeding 1 by more than 500-fold, driven by Zn (HQinh 304.7) and Al (129.5), and carcinogenic risks for Cr and As up to 9x the 1e-4 threshold. Limitation: a hazard quotient of 300 for zinc is a sign that the reference concentration is being misapplied, not that residents are inhaling a toxic zinc dose - EPA inhalation RfCs do not exist for several of these elements and the substitutes used are oral values scaled by default breathing rates; sampling duration and n are not stated. Relevance: moderate for the Central Asian composition data, which is genuinely scarce; the risk numbers should not be propagated.}
{10.3390/atmos17100977}

\paperentry{Wang et al. 2026d}
{Physics-inspired inductive spatiotemporal kriging for PM$_{2.5}$ with satellite gradient constraints}
{Atmos Environ}
{\textsc{metadata only.} Metadata-only record: Elsevier deposited title, authors and DOI on 6 Oct with no abstract. By title, this is the central methodological problem of a sparse network - interpolating to unmonitored locations - attacked inductively (so that new sensors can be added without retraining) and constrained by satellite gradients rather than by a stationarity assumption. The Everglades record in the 2 October issue measured exactly how badly ordinary kriging and IDW fail day to day; this is the proposed repair. Relevance: high, the highest of any metadata-only record in this backfill; flagged for a full-text read as soon as the abstract appears.}
{10.1016/j.atmosenv.2026.122398}

\par\vskip 0pt plus 10\baselineskip\penalty-200\vskip 0pt plus -10\baselineskip
\band{Sensing, forecasting \& instrumentation\hfill 3 records}{Amber}

\paperentry{Loh et al. 2026}
{Seasonal characteristics of submicron aerosols in the coastal atmosphere of southeastern Korea: Influence of urban-port outflow}
{Mar Pollut Bull}
{Seasonal HR-ToF-AMS measurement at a coastal receptor site southwest of Busan, winter 2021 to autumn 2022. PM$_{1}$ ranged 4.23-10.0 $\mu$g\,m$^{-3}$; PMF resolved five organic factors including hydrocarbon-like, oxidised primary (OPOA), two oxygenated and biomass-burning. Ultrafine growth events were frequent, most in autumn, when OPOA reached 36\% alongside elevated black carbon under north-easterly transport from the port. Limitation: one receptor site with no upwind reference, so 'port influence' rests on wind sector rather than on a source-resolved tracer, and AMS does not measure refractory sea salt or dust - the submicron mass is non-refractory by construction. Relevance: high - growth events at 4-10 $\mu$g\,m$^{-3}$ of mass are invisible to a mass-reporting network, which is the recurring theme of this week's issues.}
{10.1016/j.marpolbul.2026.120437}

\paperentry{Dai et al. 2026}
{Cross-phase partitioning of sulfur-nitrogen ratios and aerosol mixing-state evolution based on single-particle observations}
{Atmos Chem Phys}
{SPA-MS through winter and summer emission-control periods and matched normal periods, with three sulfur-to-nitrogen metrics (gas gSNR, particle pSNR, number-based nSNR) analysed by causal inference plus interpretable machine learning. Control periods cut NOx while SO$_2$ stayed roughly flat, raising gSNR, lifting pSNR by about 57\% across particle classes and expanding mean vacuum aerodynamic diameter by 27\%; enrichment was greatest in particles containing both black and organic carbon and least in BC-free inorganic particles. Limitation: emission-control periods are confounded with season and meteorology, and 'causal inference' on an observational time series with co-varying controls is a strong label for what remains an association. Relevance: high - a 27\% diameter shift under a control policy is a direct change in the optical cross-section a nephelometric sensor measures, i.e. a calibration drift caused by policy rather than by hardware.}
{10.5194/acp-26-14015-2026}

\paperentry{Zhang et al. 2026e}
{Full field of view temporal calibration for the directional polarimetric camera onboard the GaoFen-5 (02) satellite based on a data reconstruction method}
{Appl Opt}
{\textsc{metadata only.} Metadata-only record: PubMed indexed title, authors and DOI on 6 Oct; Applied Optics deposited no abstract and the retry leg recovered none. By title, a full-field-of-view temporal calibration for the DPC instrument, correcting on-orbit response drift by data reconstruction. Polarimetric retrievals are the satellite side of aerosol loading, and a drifting on-orbit calibration propagates straight into the AOD that surface PM$_{2.5}$ models are trained against - the same class of problem as field drift in a ground sensor, at a different scale. Relevance: moderate-high; on the retry list.}
{10.1364/ao.597069}

\par\vskip 0pt plus 10\baselineskip\penalty-200\vskip 0pt plus -10\baselineskip
\band{Occupational \& indoor\hfill 1 record}{Amber}

\paperentry{Othman et al. 2026}
{Impact of Urban Land Use on Indoor Air Quality: Quantifying Ambient Particle Infiltration and Asthma Risk in Naturally Ventilated School Buildings}
{Indoor Air}
{Ten primary schools stratified by land use (high-traffic, industrial, residential, mixed industrial/high-traffic, mixed industrial/port) with paired indoor-outdoor PM$_{2.5}$ and PM$_{1}$. Indoor maxima were at mixed industrial/port schools (24 h PM$_{2.5}$ 35.33 +/- 9.25, PM$_{1}$ 23.79 +/- 5.87 $\mu$g\,m$^{-3}$); outdoor maxima at high-traffic schools (34.50 +/- 20.64 and 23.78 +/- 13.09); land-use class differences were significant (p<0.05) and the mean infiltration factor exceeded 0.6. Limitation: ten buildings across five classes is two per class, so the land-use contrast is confounded with building; Finf estimated from indoor-outdoor ratios without an independent air-change measurement absorbs any indoor source into the infiltration term. Relevance: high - Finf > 0.6 in naturally ventilated classrooms means outdoor network data is a usable proxy for children's exposure in this climate, which is the assumption most school studies make without testing.}
{10.1155/ina/5050009}

\par\vskip 0pt plus 10\baselineskip\penalty-200\vskip 0pt plus -10\baselineskip
\band{Burden, policy \& mitigation\hfill 2 records}{Coral}

\paperentry{Ranzi et al. 2026}
{Atlante Aria e Salute: methods for estimating the health impacts of air pollution and national application}
{Epidemiol Prev}
{Methodological description of the Air and Health Atlas, an interactive national platform. PM$_{2.5}$ and NO$_2$ are modelled at 1 km with random forests, converted to population-weighted exposure at municipality level against 2021 census blocks, and combined with provincial death-certificate rates through HRAPIE-2 concentration-response functions against the 5 $\mu$g\,m$^{-3}$ WHO counterfactual; 2019 is shown as an illustration, with pronounced north-south gradients. Limitation: population-weighted exposure at municipality scale cannot capture within-municipality contrast, the random-forest surface is validated against the same monitors that train it, and HRAPIE-2 functions are transferred from other populations. Relevance: high - this is the operational template a low-cost network would be asked to improve, and the 1 km population-weighted step is exactly where added density pays.}
{10.19191/ep26.4-5.s1.a1030.092}

\paperentry{Zhao et al. 2026e}
{Thermal-oxygen evolution and atmospheric pollutant emission characteristics during the vertical propagation of underground coal fires}
{Research Square (preprint)}
{A scaled physical-similarity rig run from ignition to burnout with continuous monitoring at six depths. Thermal-oxygen evolution split into incubation, rapid heating, intensive reaction, rapid cooling and burnout; the upper layers burned out while the middle sustained reaction and the deepest was oxygen-limited; surface response lagged the seam. Particulate matter and light hydrocarbons were released early (CH4 and C2H6 concentrated in the first 40-60 h, at background by \textasciitilde{}80 h) while CO$_2$ and CO dominated later. Limitation: a similarity-scaled rig reproduces sequence, not magnitude - emission factors from it cannot be transferred to a real coal fire, and no scaling validation is offered; preprint. Relevance: moderate - the early-PM, late-CO sequence is a usable detection signature for surface sensor arrays over suspected underground fires.}
{10.21203/rs.3.rs-10517753/v1}

\clearpage
\section*{\sffamily\bfseries\color{Deep}Corpus register}
\vspace{-1.5mm}{\color{Amber}\rule{\linewidth}{1.1pt}}\vspace{2.5mm}

\input{table_rows}

\vspace{2mm}

\begin{keybox}
{\sffamily\bfseries\footnotesize\color{Deep}Provenance and method}\\[1.2mm]
{\fontsize{7.2}{8.8}\selectfont
Entry window \textbf{6 October 2026}, single day, continuous with the 5 October issue. Harvest
and screening were completed on the \textbf{6 October} run and cached
(\texttt{cache/2026-10-06/\_fresh.json}); this run authored from that cache without re-querying.
\textbf{72} fresh candidates after cross-day deduplication, including \textbf{7} appended from the PubMed
connector; \textbf{OpenAlex} HTTP 429 again; \textbf{arXiv} answered (148\,kB) with no entry dated
6 October passing the recency gate. \textbf{11 admitted}, \textbf{61 rejected} with logged reasons.
\textbf{Four records carried on metadata alone} (36\,\%) --- the fourth consecutive issue above 25\,\%,
and one of them (\textit{Atmos.\ Environ.}, satellite-constrained inductive kriging) is the
highest-relevance record of the day despite having no abstract. \textbf{Consensus} returned 10 calibration
hits, \textbf{0 new}. \textbf{Scholar Gateway} fails with \texttt{ACCESS\_DENIED}.
\textbf{ClinicalTrials.gov} (LastUpdatePostDate 3--6 October): EPIC-AIR \texttt{NCT07500948} and a
CAN-EXACT update, both already in the trial watch. An \textit{Atmos.\ Environ.} retraction notice
(drought and PM$_{2.5}$ in Texas) appeared in the window and concerns no record in this corpus. DOI gate:
\textbf{0 fail, 1 warn} --- the \textit{Epidemiologia \& Prevenzione} supplement DOI resolves but is not in
Crossref, which is normal for that publisher and was confirmed by eye. Abstracts remain the copyright of
their respective publishers.\par}
\end{keybox}

\end{document}
